\FloatBarrier
\section{本章小结}
\label{sec:multivariate-differentiation-chapter-summary}

本章的对象是位置上的量：标量场在每点给出一个数，向量场在每点给出一个与位置同维的向量，一般向量值函数的输出维数可以不同。它们的极限与连续性按输入、输出的范数理解，输出可以逐分量判断，输入的多个坐标却不能拆开逐个检查；沿两条路径得到不同极限能否定极限存在，只查若干条路径则不能证明它存在。

偏导数和方向导数只考察固定的直线。可微要求所有小扰动共用一个线性近似，包括方向随步长改变的扰动；这个近似由偏导数唯一确定。偏导数存在、全部方向导数存在或函数连续，都不单独保证可微；即使全部方向导数组成一个线性映射，仍须满足统一的余项条件。偏导数连续是方便的充分条件，但不是必要条件。

全微分是输入扰动上的线性泛函，梯度是它在所选内积下的表示向量。改变内积会改变梯度的坐标和最陡上升方向，全微分本身不变。沿含于定义域的线段，中值公式把梯度的大小化为有限变化的界；在凸开集上，梯度长度的统一上界恰好就是Lipschitz常数，凸性不能省略。

输出为向量时，微分由Jacobian矩阵表示，它的行是各分量的梯度，列是各输入坐标对输出点的推动。复合映射的Jacobian按计算顺序相乘，标量目标的梯度则由微分的伴随传回输入端。向量增量没有共同的中值点，改用沿线段的积分表示，并由导数的算子范数控制长度。

把微分再用于梯度，得到Hessian与二阶Taylor公式。混合偏导数在连续时可以交换，$C^2$函数的Hessian因而对称。在开集内的驻点处，Hessian正定、负定或不定分别对应严格极小、严格极大和鞍点；半定而退化的情形需要更高阶的信息，边界点和受约束的问题也不在这一判别的范围内。

这些结论都是局部的：可微只保证余项是$o(\lVert\bm h\rVert)$，不说明给定步长下误差有多大。矩阵自变量的微分公式、隐式方程确定的解，以及导数怎样控制误差放大和一阶预测的精度，依次在后面三章展开。
